\documentclass[11pt]{article}
\usepackage[margin=1in]{geometry}
\usepackage[T1]{fontenc}
\usepackage[utf8]{inputenc}
\usepackage{lmodern}
\usepackage{microtype}
\usepackage[table]{xcolor}
\usepackage{amsmath,amssymb,mathtools}
\usepackage{booktabs,array,longtable,tabularx}
\usepackage{hyperref}
\usepackage{enumitem}
\usepackage{fancyhdr}
\usepackage{titlesec}
\usepackage{tcolorbox}
\hypersetup{colorlinks=true, linkcolor=blue!50!black, urlcolor=blue!50!black}

\definecolor{TitleBlue}{HTML}{163A5F}
\definecolor{AccentTeal}{HTML}{2D6F73}
\definecolor{SoftBlue}{HTML}{EAF3F8}
\definecolor{SoftTeal}{HTML}{EDF8F6}
\definecolor{SoftGray}{HTML}{F5F7FA}

\pagestyle{fancy}
\fancyhf{}
\lhead{PINN for Two-Dimensional MHD Evolution}
\rhead{Implementation Note}
\cfoot{\thepage}
\setlength{\headheight}{14pt}

\titleformat{\section}{\Large\bfseries\color{TitleBlue}}{\thesection}{0.6em}{}
\titleformat{\subsection}{\large\bfseries\color{AccentTeal}}{\thesubsection}{0.6em}{}

\newtcolorbox{overviewbox}{colback=SoftBlue,colframe=TitleBlue,boxrule=0.8pt,arc=2mm,left=3mm,right=3mm,top=2mm,bottom=2mm}
\newtcolorbox{implbox}{colback=SoftTeal,colframe=AccentTeal,boxrule=0.8pt,arc=2mm,left=3mm,right=3mm,top=2mm,bottom=2mm}

\begin{document}

\begin{center}
    {\LARGE \bfseries \color{TitleBlue} Implementation Note: MATLAB File Map and Dependency Guide}\\[0.5em]
    {\large \color{AccentTeal} Physics-Informed Neural Network for Two-Dimensional MHD Evolution}
\end{center}

\vspace{0.5em}

\begin{overviewbox}
\textbf{Purpose.} This note explains what each major MATLAB file does, how the folders are organized, and how the pieces depend on one another. It is intended as a practical orientation document for someone trying to understand or present the repository quickly.
\end{overviewbox}

\section{Top-level files}
\begin{center}
\begin{tabular}{p{0.28\linewidth} p{0.60\linewidth}}
\toprule
\textbf{File} & \textbf{Role} \\
\midrule
\texttt{README.md} & top-level project description and repository entry point \\
\texttt{project\_plan.md} & staged development plan and project roadmap \\
\texttt{run\_all.m} & intended umbrella entry point for broad execution workflows \\
\bottomrule
\end{tabular}
\end{center}

\section{Folder overview}
The repository separates code by function:
\begin{itemize}[leftmargin=1.8em]
    \item \texttt{src/physics}: physical variable conversions, thermodynamic helpers, physical fluxes, and state-sanity utilities,
    \item \texttt{src/numerics}: boundary conditions, reconstruction, numerical flux assembly, right-hand-side construction, and time integrators,
    \item \texttt{src/init}: initial-condition builders,
    \item \texttt{src/diagnostics}: error norms, phase-speed estimates, conservation metrics, and magnetic-divergence diagnostics,
    \item \texttt{src/plotting}: plotting utilities for verification and nonlinear demonstration cases,
    \item \texttt{src/grid}: grid generation,
    \item \texttt{src/utils}: small reusable helper functions,
    \item \texttt{tests}: low-level and regression-style test scripts,
    \item \texttt{cases}: reproducible top-level solver runs,
    \item \texttt{pinn}: physics-informed neural network routines,
    \item \texttt{docs}: physics notes and implementation notes.
\end{itemize}

\section{Source files: physics layer}
\begin{center}
\begin{tabular}{p{0.33\linewidth} p{0.55\linewidth}}
\toprule
\textbf{File} & \textbf{Role} \\
\midrule
\texttt{primitive\_to\_conserved.m} & maps primitive variables $(\rho,u_x,u_y,u_z,p,B_x,B_y,B_z)$ to conserved variables \\
\texttt{conserved\_to\_primitive.m} & recovers primitive variables from the conserved state \\
\texttt{compute\_pressure.m} & reconstructs gas pressure from conserved variables \\
\texttt{compute\_total\_energy.m} & builds total energy density from primitive variables \\
\texttt{flux\_x.m} & physical ideal-MHD flux in the $x$ direction \\
\texttt{flux\_y.m} & physical ideal-MHD flux in the $y$ direction \\
\texttt{fast\_magnetosonic\_speed\_x.m} & directional fast-wave speed estimate for interfaces normal to $x$ \\
\texttt{fast\_magnetosonic\_speed\_y.m} & directional fast-wave speed estimate for interfaces normal to $y$ \\
\texttt{default\_positivity\_settings.m} & creates default positivity-floor configuration \\
\texttt{apply\_positivity\_floors.m} & enforces density and pressure floors after an update when enabled \\
\texttt{check\_physical\_state.m} & sanity checks for physically admissible states \\
\bottomrule
\end{tabular}
\end{center}

\section{Source files: numerics layer}
\begin{center}
\begin{tabular}{p{0.33\linewidth} p{0.55\linewidth}}
\toprule
\textbf{File} & \textbf{Role} \\
\midrule
\texttt{apply\_periodic\_bc.m} & adds one layer of periodic ghost cells around the physical domain \\
\texttt{compute\_time\_step.m} & computes the explicit time step from the Courant-Friedrichs-Lewy condition \\
\texttt{rusanov\_flux\_x.m} & numerical interface flux in the $x$ direction \\
\texttt{rusanov\_flux\_y.m} & numerical interface flux in the $y$ direction \\
\texttt{compute\_rhs.m} & first-order semi-discrete right-hand side using piecewise-constant states \\
\texttt{compute\_rhs\_plm.m} & higher-order right-hand side using piecewise-linear reconstruction \\
\texttt{reconstruct\_plm\_minmod.m} & piecewise-linear state reconstruction with the minmod limiter \\
\texttt{update\_fv\_euler.m} & forward-Euler finite-volume update \\
\texttt{update\_fv\_rk2.m} & second-order Runge-Kutta update for the first-order spatial operator \\
\texttt{update\_fv\_rk2\_plm.m} & second-order Runge-Kutta update paired with piecewise-linear reconstruction \\
\bottomrule
\end{tabular}
\end{center}

\begin{implbox}
\textbf{Dependency idea.} The numerics layer calls the physics layer repeatedly. For example, a time-step routine needs wave speeds, and a Rusanov flux routine needs both physical fluxes and directional characteristic speeds.
\end{implbox}

\section{Source files: initialization, diagnostics, plotting, and utilities}
\subsection{Initialization}
\begin{center}
\begin{tabular}{p{0.33\linewidth} p{0.55\linewidth}}
\toprule
\textbf{File} & \textbf{Role} \\
\midrule
\texttt{init\_alfven\_wave.m} & builds the small-amplitude shear-Alfv\'en verification initial condition \\
\texttt{init\_magnetic\_pressure\_pulse.m} & builds the nonlinear magnetic-pressure-pulse initial condition \\
\bottomrule
\end{tabular}
\end{center}

\subsection{Diagnostics}
\begin{center}
\begin{tabular}{p{0.33\linewidth} p{0.55\linewidth}}
\toprule
\textbf{File} & \textbf{Role} \\
\midrule
\texttt{compute\_l2\_error.m} & computes a discrete $L^2$ error norm \\
\texttt{measure\_alfven\_phase\_speed.m} & estimates the Alfv\'en phase speed from periodic line-cut data \\
\texttt{compute\_divB.m} & estimates magnetic divergence on the grid \\
\texttt{compute\_total\_mass.m} & computes total mass over the domain \\
\texttt{compute\_total\_energy\_domain.m} & computes total energy over the domain \\
\bottomrule
\end{tabular}
\end{center}

\subsection{Plotting and utilities}
\begin{center}
\begin{tabular}{p{0.33\linewidth} p{0.55\linewidth}}
\toprule
\textbf{File} & \textbf{Role} \\
\midrule
\texttt{plot\_alfven\_linecuts.m} & numerical-versus-exact line cuts for the Alfv\'en benchmark \\
\texttt{plot\_pulse\_fields.m} & field plots and visualization for the magnetic-pressure-pulse case \\
\texttt{make\_uniform\_grid.m} & creates the uniform Cartesian grid object used throughout the project \\
\texttt{minmod.m} & limiter helper used by the piecewise-linear reconstruction routine \\
\bottomrule
\end{tabular}
\end{center}

\section{Case scripts}
\begin{center}
\begin{tabular}{p{0.33\linewidth} p{0.55\linewidth}}
\toprule
\textbf{File} & \textbf{Role} \\
\midrule
\texttt{alfven\_wave\_case.m} & primary verification run for the small-amplitude Alfv\'en benchmark \\
\texttt{alfven\_wave\_convergence\_study.m} & repeated verification run across multiple grid sizes \\
\texttt{magnetic\_pressure\_pulse\_case.m} & nonlinear demonstration problem for localized magnetic overpressure \\
\bottomrule
\end{tabular}
\end{center}

The case scripts should stay thin. Their job is to assemble a complete experiment from reusable lower-level components rather than to contain duplicated numerical logic.

\section{Test scripts}
\begin{center}
\begin{tabular}{p{0.33\linewidth} p{0.55\linewidth}}
\toprule
\textbf{File} & \textbf{Role} \\
\midrule
\texttt{test\_variable\_conversion.m} & variable conversion consistency check \\
\texttt{test\_flux\_functions.m} & ideal-MHD physical flux sanity checks \\
\texttt{test\_time\_step\_and\_wave\_speeds.m} & time-step and characteristic-speed checks \\
\texttt{test\_uniform\_state\_preservation.m} & first-order uniform-state preservation \\
\texttt{test\_rk2\_uniform\_state\_preservation.m} & Runge-Kutta uniform-state preservation \\
\texttt{test\_plm\_uniform\_state\_preservation.m} & piecewise-linear uniform-state preservation \\
\texttt{test\_positivity\_floor\_application.m} & positivity-floor behavior check \\
\bottomrule
\end{tabular}
\end{center}

\section{PINN files}
\subsection{Current PINN philosophy}
The PINN side is split into two layers:
\begin{enumerate}[leftmargin=1.8em]
    \item a working reduced shear-Alfv\'en PINN benchmark,
    \item a broader full ideal-MHD residual scaffold that is not yet a complete validated training pipeline.
\end{enumerate}

\begin{center}
\begin{tabular}{p{0.33\linewidth} p{0.55\linewidth}}
\toprule
\textbf{File} & \textbf{Role} \\
\midrule
\texttt{pinn\_main\_alfven.m} & top-level entry point for the reduced Alfv\'en PINN benchmark \\
\texttt{pinn\_default\_params.m} & default network, sampling, and optimizer parameters \\
\texttt{pinn\_build\_network.m} & neural-network construction \\
\texttt{pinn\_sample\_collocation\_points.m} & collocation-point generation for residual enforcement \\
\texttt{pinn\_predict\_fields.m} & network output interpretation as physical fields \\
\texttt{pinn\_apply\_ic\_bc.m} & initial-condition and boundary-condition loss terms \\
\texttt{pinn\_ideal\_mhd\_residuals\_alfven.m} & reduced shear-Alfv\'en residuals derived from ideal MHD \\
\texttt{pinn\_model\_gradients\_alfven.m} & gradient assembly for reduced PINN training \\
\texttt{pinn\_compute\_losses.m} & benchmark loss breakdown for the reduced Alfv\'en PINN \\
\texttt{pinn\_train\_alfven.m} & reduced PINN training loop \\
\texttt{pinn\_plot\_training\_history.m} & training-history visualization \\
\texttt{pinn\_plot\_comparison\_alfven.m} & solver-versus-PINN comparison plots \\
\texttt{pinn\_compare\_with\_solver.m} & comparison wrapper against the conventional solver \\
\texttt{pinn\_ideal\_mhd\_residuals\_general.m} & scaffold for broader 2D ideal-MHD residual evaluation \\
\texttt{pinn\_compute\_losses\_general.m} & scaffolded loss wrapper for the broader residual engine \\
\bottomrule
\end{tabular}
\end{center}

\section{Dependency ladder for understanding the code}
The most efficient way to understand the repository is:
\begin{enumerate}[leftmargin=1.8em]
    \item start with the state-conversion and flux functions,
    \item move to the numerical flux and right-hand-side assembly,
    \item then inspect the update routines,
    \item then read the initialization and diagnostics files,
    \item then look at the case scripts,
    \item only after that inspect the PINN folder.
\end{enumerate}

\section{What is stable versus still exploratory}
\textbf{Relatively stable at the current stage:}
\begin{itemize}[leftmargin=1.8em]
    \item grid generation,
    \item state conversion,
    \item physical flux routines,
    \item Rusanov-based finite-volume solver path,
    \item Alfv\'en-wave and pulse-case script structure,
    \item low-level tests and preservation tests.
\end{itemize}

\textbf{Still exploratory or scaffolded:}
\begin{itemize}[leftmargin=1.8em]
    \item broader full ideal-MHD PINN training pipeline,
    \item general residual normalization and loss balancing,
    \item more advanced benchmarks such as a truly oblique 2D Alfv\'en wave.
\end{itemize}

\begin{overviewbox}
\textbf{Bottom line.} The repository is easiest to understand if it is read as a ladder: physics layer $\rightarrow$ numerics layer $\rightarrow$ diagnostics and cases $\rightarrow$ PINN benchmark $\rightarrow$ general PINN scaffold. The conventional solver is already the core of the project; the PINN side is built on top of that foundation.
\end{overviewbox}

\end{document}
